\label{wb:lem:toric-chart-control}
For a torus point \(p=(x,t)\) with \(0<|t|<1\), put
\[
 y(p)=\frac{(\log|x_1|,\log|x_2|)}{\log|t|}.
\]
Then
\[
 p\in\mathcal U_\sigma(S)
 \quad\Longleftrightarrow\quad
 y(p)\in\sigma^{(\log S/|\log|t||)}.
 \tag{2.5c}
\]
The exact closed-unit-polydisc version is
\begin{equation}
 p\in\overline{\mathcal U}_\sigma(1)
 \quad\Longleftrightarrow\quad
 \theta_i^\sigma(y(p))\ge0\ (i=0,1,2)
 \quad\Longleftrightarrow\quad y(p)\in\sigma.
\label{wb:eq:cusp-closed-chart-identity}
\end{equation}
Every point of \(\{|t|<1\}\subset Y\) belongs to
\(\overline{\mathcal U}_\sigma(1)=\{|z_i|\le1\}\) for some \(\sigma\).
Consequently the sets \(\mathcal U_\sigma(2,\epsilon)\) cover \(Y_\epsilon\), and
every compact subset is contained in the union of finitely many members of
this cover.
\end{lemma}
\Uses{the hypotheses, notation and definitions stated here; the explicit constructions and calculations in the proof}


\begin{proof}
Since \((y,1)=\sum_i\theta_i^\sigma(y)(v_i,1)\), duality gives
\[
 \log|z_i(p)|=\log|t(p)|\,\theta_i^\sigma(y(p)).
\]
Because \(\log|t|<0\), this is exactly (2.5c); retaining non-strict
inequalities proves the closed-unit-polydisc assertion, including triangle
edges and vertices.

For the boundary assertion let \(p\) lie in the toric orbit
\(O_{\widehat P}\) of the cone over a cell \(P\).  In a chart with
\(\sigma\supseteq P\), the coordinates corresponding to vertices of
\(P\) vanish.  On the quotient lattice
\(N'/(N'\cap\operatorname{span}\widehat P)\), the remaining coordinates
are the characters dual to the image cone \(\bar\sigma\).  Hence
\(p\in\overline{\mathcal U}_\sigma(1)\) exactly when
\(-\operatorname{Log}_{\widehat P}(p)\in\bar\sigma\).
The cones \(\bar\sigma\), for triangles \(\sigma\supseteq P\), cover the
quotient real vector space: lift a vector \(n\), choose a relative interior
point \(p_0\in P\times\{1\}\), and for sufficiently small \(u>0\) place
the height-one rescaling of \(p_0+un\) in a triangle of the locally finite
triangulation.  That triangle contains \(P\), and its quotient cone contains
\(n\).  This proves the asserted cover.  Compactness then supplies a finite
subcover.
\end{proof}